% -*- coding: utf-8 -*-
\documentclass[12pt,a4paper]{article}
\usepackage[T2A]{fontenc}
\usepackage[utf8]{inputenc}
\usepackage[english,russian]{babel}
\usepackage{amsmath,amssymb,amsthm}
\usepackage{geometry}
\usepackage{graphicx}
\usepackage{hyperref}
\usepackage{booktabs}
\usepackage{caption}
\geometry{left=25mm,right=25mm,top=25mm,bottom=25mm}

\title{Домашнее задание — Блок 2\newline Отчёт}
\date{}

\begin{document}
\maketitle
\tableofcontents
\clearpage

\section{Постановка задачи}
Рассматривается функция двух переменных
\[
z(x,y)=x^{3}+3x^{2}+y^{3}-5y^{2}+7y-7.
\]
Требуется:
\begin{itemize}
  \item визуализировать линии уровня и траектории итераций (ДЗ2.1),
  \item найти точку минимума численно методом Ньютона и выполнить аналитическое исследование стационарных точек (ДЗ2.2),
  \item представить объектную модель для методов оптимизации (ДЗ2.3).
\end{itemize}

\section{Краткие численные результаты}
Ниже перечислены результаты, полученные ранее и собранные в текстовом отчёте:
\begin{itemize}
  \item Метод покоординатного спуска: найдено около $x^*=(-0.000013,\,2.333317)$, $z^*\approx -5.185185$, число итераций: 2.
  \item Метод градиентного спуска: $x^*\approx(-0.003558,\,2.338417)$, $z^*\approx -5.185095$, итераций: 15.
  \item Метод наискорейшего спуска: $x^*\approx(0.000009,\,2.333330)$, $z^*\approx -5.185185$, итераций: 1.
  \item Метод Ньютона: $x^*=(0.000000,\,2.333333)$, $z^*\approx -5.185185$, итераций: 4.
\end{itemize}

\subsection{Формулы методов оптимизации}
Для этой же задачи используются три классических метода спуска. Ниже записаны их рабочие формулы.

\paragraph{Покоординатный спуск.}
На каждой итерации последовательно минимизируются одномерные функции
\[
\varphi_1(t)=z(t,y^{(k)}),\qquad \varphi_2(t)=z(x^{(k+1)},t).
\]
Обновление имеет вид
\[
x^{(k+1)}=\arg\min_{x\in[x^{(k)}-s,\,x^{(k)}+s]} z(x,y^{(k)}),
\qquad
y^{(k+1)}=\arg\min_{y\in[y^{(k)}-s,\,y^{(k)}+s]} z(x^{(k+1)},y),
\]
где $s>0$ — локальный радиус поиска.

\paragraph{Градиентный спуск.}
Используется формула
\[
\begin{pmatrix}
x^{(k+1)}\\
y^{(k+1)}
\end{pmatrix}
=
\begin{pmatrix}
x^{(k)}\\
y^{(k)}
\end{pmatrix}
-\alpha_k
\nabla z(x^{(k)},y^{(k)}),
\]
где
\[
\nabla z(x,y)=
\begin{pmatrix}
3x^2+6x\\
3y^2-10y+7
\end{pmatrix}.
\]
В реализации шаг $\alpha_k$ уменьшается вдвое, если новая точка не дает убывания функции.

\paragraph{Метод наискорейшего спуска.}
На каждом шаге направление берется антиградиентом,
\[
d_k=-\nabla z(x^{(k)},y^{(k)}),
\]
а шаг $h_k$ выбирается из одномерной задачи
\[
h_k=\arg\min_{h\in[h_c-s_h,\,h_c+s_h]} z\bigl(x^{(k)}-h\,z_x(x^{(k)},y^{(k)}),\,y^{(k)}-h\,z_y(x^{(k)},y^{(k)})\bigr),
\]
где $h_c$ — текущий центр поиска, а $s_h$ — локальный радиус поиска.

\section{Аналитическое решение (ДЗ2.2)}
\subsection{Найдем стационарные точки}
Стационарные точки решениями системы
\[
\frac{\partial z}{\partial x}=0,\qquad \frac{\partial z}{\partial y}=0.
\]
Вычислим частные производные:
\[
\frac{\partial z}{\partial x}=3x^{2}+6x=3x(x+2),
\]
\[
\frac{\partial z}{\partial y}=3y^{2}-10y+7.
\]
Решаем по очереди:
\begin{itemize}
  \item из $3x(x+2)=0$ получаем $x=0$ или $x=-2$.
  \item квадратное уравнение для $y$: $3y^{2}-10y+7=0$. Дискриминант
  \[D=(-10)^{2}-4\cdot3\cdot7=100-84=16.
  \]
  Тогда
  \[y=\frac{10\pm\sqrt{16}}{6}=\frac{10\pm4}{6}.
  \]
  Получаем $y=1$ и $y=7/3$.
\end{itemize}
Следовательно, все стационарные точки (все комбинации) — это четыре точки:
\[
(0,1),\quad (0,7/3),\quad (-2,1),\quad (-2,7/3).
\]

\subsection{Классификация по матрице Гессе}
Вычислим элементы матрицы Гессе $H=\begin{pmatrix} z_{xx} & z_{xy}\\ z_{yx} & z_{yy}\end{pmatrix}$:
\[
z_{xx}=6x+6,\qquad z_{xy}=z_{yx}=0,\qquad z_{yy}=6y-10.
\]
Определитель $\det H=(6x+6)(6y-10)$ и след $\operatorname{tr}H=6x+6+6y-10$.

Рассмотрим каждую точку:
\begin{itemize}
  \item $(0,1)$: $z_{xx}=6$, $z_{yy}=-4$, $\det H=6\cdot(-4)=-24<0$ \textit{седловая точка}.
  \item $(0,7/3)$: $z_{xx}=6$, $z_{yy}=6\cdot\frac{7}{3}-10=4$, $\det H=6\cdot4=24>0$, $z_{xx}>0$ \textit{локальный минимум}.
  \item $(-2,1)$: $z_{xx}=6(-2)+6=-6$, $z_{yy}=-4$, $\det H=(-6)(-4)=24>0$, $z_{xx}<0$ \textit{локальный максимум}.
  \item $(-2,7/3)$: $z_{xx}=-6$, $z_{yy}=4$, $\det H=-6\cdot4=-24<0$ \textit{седловая точка}.
\end{itemize}

\subsection{Значения функции в стационарных точках}
Подставим координаты в $z(x,y)$:
\begin{align*}
z(0,1)&=0+0+1-5+7-7=-4,\\
z\bigl(0,\tfrac{7}{3}\bigr)&=\left(\tfrac{7}{3}\right)^{3}-5\left(\tfrac{7}{3}\right)^{2}+7\left(\tfrac{7}{3}\right)-7=-5.185185\ldots,\\
z(-2,1)&=(-8)+3\cdot4+1-5+7-7=0,\\
z\bigl(-2,\tfrac{7}{3}\bigr)&=4-5.185185\ldots=-1.185185\ldots.
\end{align*}

Из анализа видно, что глобального (по всей $R^{2}$) минимума у этой функции нет (кубические члены), но среди стационарных точек локальный минимум достигается в точке $(0,7/3)$ со значением $\approx -5.185185$.

\section{Численный метод Ньютона (ДЗ2.2)}
\subsection{Метод и критерий остановки}
Численный метод Ньютона для задачи минимизации сводится к решению системы стационарности
\[
\nabla z(x,y)=0.
\]
В одномерном случае метод Ньютона имеет вид
\[
x_{k+1}=x_k-\frac{f'(x_k)}{f''(x_k)}.
\]
В двумерном случае на каждой итерации решается линейная система
\[
H(x^{(k)})\,\delta^{(k)}=\nabla z(x^{(k)}),\qquad x^{(k+1)}=x^{(k)}-\delta^{(k)}.
\]
Критерий останова используется по норме градиента: $\|\nabla z(x^{(k)})\|<10^{-4}$.

\subsection{Подготовка формул для данной функции}
Для функции
\[
z(x,y)=x^3+3x^2+y^3-5y^2+7y-7
\]
частные производные и матрица Гессе равны
\[
\frac{\partial z}{\partial x}=3x^2+6x,\qquad
\frac{\partial z}{\partial y}=3y^2-10y+7,
\]
\[
H(x,y)=
\begin{pmatrix}
6x+6 & 0\\
0 & 6y-10
\end{pmatrix}.
\]
Поскольку матрица Гессе диагональна, система Ньютона распадается на две независимые одномерные формулы:
\[
x_{k+1}=x_k-\frac{3x_k^2+6x_k}{6x_k+6}
=\frac{x_k^2}{2(x_k+1)},
\]
\[
y_{k+1}=y_k-\frac{3y_k^2-10y_k+7}{6y_k-10}
=\frac{3y_k^2-7}{6y_k-10}.
\]

\subsection{Ручное решение по итерациям}
В качестве начальной точки берём точку из условия:
\[
(x_0,y_0)=\left(-\frac12,\,\frac52\right).
\]
Далее выполняем ручной расчёт каждой итерации.

\paragraph{Итерация 0.}
\[
\nabla z\left(-\frac12,\frac52\right)=\left(-\frac94,\,\frac34\right),
\qquad
H\left(-\frac12,\frac52\right)=
\begin{pmatrix}
3 & 0\\
0 & 5
\end{pmatrix}.
\]
Тогда поправка Ньютона равна
\[
\delta^{(0)}=H^{-1}\nabla z=
\left(-\frac34,\,\frac3{20}\right).
\]
Следующая точка:
\[
x_1=x_0-\delta_x^{(0)}=-\frac12-\left(-\frac34\right)=\frac14,
\qquad
y_1=y_0-\delta_y^{(0)}=\frac52-\frac3{20}=\frac{47}{20}.
\]

\paragraph{Итерация 1.}
\[
\nabla z\left(\frac14,\frac{47}{20}\right)=\left(\frac{27}{16},\,\frac{27}{400}\right),
\qquad
H\left(\frac14,\frac{47}{20}\right)=
\begin{pmatrix}
\frac{15}{2} & 0\\
0 & \frac{41}{10}
\end{pmatrix}.
\]
Поправка:
\[
\delta^{(1)}=
\left(\frac{9}{40},\,\frac{27}{1640}\right).
\]
Следующая точка:
\[
x_2=\frac14-\frac{9}{40}=\frac1{40},
\qquad
y_2=\frac{47}{20}-\frac{27}{1640}=\frac{3827}{1640}.
\]

\paragraph{Итерация 2.}
\[
\nabla z\left(\frac1{40},\frac{3827}{1640}\right)=
\left(\frac{243}{1600},\,\frac{2187}{2689600}\right),
\]
\[
H\left(\frac1{40},\frac{3827}{1640}\right)=
\begin{pmatrix}
\frac{123}{20} & 0\\
0 & \frac{3281}{820}
\end{pmatrix}.
\]
Поправка:
\[
\delta^{(2)}=
\left(\frac{81}{3280},\,\frac{2187}{10761680}\right).
\]
Следующая точка:
\[
x_3=\frac1{40}-\frac{81}{3280}=\frac1{3280},
\qquad
y_3=\frac{3827}{1640}-\frac{2187}{10761680}=\frac{25110587}{10761680}.
\]

\paragraph{Итерация 3.}
\[
\nabla z\left(\frac1{3280},\frac{25110587}{10761680}\right)=
\left(\frac{19683}{10758400},\,\frac{14348907}{115813756422400}\right),
\]
\[
H\left(\frac1{3280},\frac{25110587}{10761680}\right)=
\begin{pmatrix}
\frac{9843}{1640} & 0\\
0 & \frac{21523361}{5380840}
\end{pmatrix}.
\]
Поправка:
\[
\delta^{(3)}=
\left(\frac{6561}{21523360},\,\frac{14348907}{463255047212960}\right).
\]
Следующая точка:
\[
x_4=\frac1{3280}-\frac{6561}{21523360}=\frac1{21523360},
\qquad
y_4=\frac{25110587}{10761680}-\frac{14348907}{463255047212960}
\approx 2.3333333333.
\]

\paragraph{Проверка критерия остановки.}
В точке $(x_4,y_4)$ получаем
\[
\|\nabla z(x_4,y_4)\|\approx 2.79\cdot 10^{-7}<10^{-4},
\]
значит метод останавливается на 4-й итерации. Итоговая точка совпадает с локальным минимумом
\[
(x^*,y^*)=\left(0,\,\frac73\right),
\qquad
z(x^*,y^*)=-5.185185\ldots
\]

\subsection{Сводка численного результата}
Выполнен запуск численного метода Ньютона. Итог:
\begin{itemize}
  \item найденная точка: $x^*=(0,\,7/3)$,
  \item значение функции: $z(x^*)\approx -5.185185$,
  \item число итераций: 4.
\end{itemize}

\section{Объектная модель}
Объектная модель решения оформлена вокруг абстракции оптимизатора и общей структуры истории итераций.

\subsection{Состав модели}
В файле \texttt{model/base.py} определены:
\begin{itemize}
  \item \texttt{Optimizer2D} — абстрактный базовый класс с методом \texttt{run()},
  \item \texttt{IterationPoint} — одна запись истории итераций,
  \item \texttt{OptimizationResult} — итог работы метода оптимизации.
\end{itemize}

В файле \texttt{model/methods.py} реализованы конкретные стратегии:
\begin{itemize}
  \item \texttt{CoordinateDescentOptimizer},
  \item \texttt{GradientDescentOptimizer},
  \item \texttt{SteepestDescentOptimizer},
  \item \texttt{NewtonOptimizer2D}.
\end{itemize}

\subsection{Связи между классами}
Используются следующие отношения объектной модели:
\begin{itemize}
  \item \textbf{Наследование:} все численные методы наследуются от \texttt{Optimizer2D} и реализуют единый интерфейс \texttt{run()}.
  \item \textbf{Композиция:} объект \texttt{OptimizationResult} содержит список \texttt{IterationPoint}, то есть полную историю вычислений.
  \item \textbf{Использование:} визуализация и генерация отчёта работают с результатами методов, не зная внутренней реализации алгоритмов.
\end{itemize}

\subsection{Поток работы программы}
Сначала создается конфигурация задачи, затем последовательно запускаются все методы оптимизации. После этого:
\begin{enumerate}
  \item строятся графики линий уровня и траекторий,
  \item формируется аналитический вывод по стационарным точкам,
  \item результаты собираются в итоговый отчет.
\end{enumerate}

Такое разделение удобно тем, что алгоритмы, визуализация и оформление отчета отделены друг от друга и могут развиваться независимо.
\end{document}
